\section{Additional Implementation Details}

In this section, we describe some additional implementation details of \ourmodel{} ---
its architecture and training ---
and implementation details for our off-the-shelf 3D detectors.

\subsection{Architecture}

\paragraph{Trajectory Forecasting Header:}


% Inspired by \cite{casas2024detra} and \cite{chitta2022transfuser},
Each refinement block predicts trajectory forecasts from the refinement features $\mathbf{Q}^{\mathrm{ref}(i)} \in \mathbb{R}^{N^\mathrm{ref} \times T_f + 1 \times d}$ using a
single-layer bi-directional GRU~\cite{cho2014learning}.
The output of the GRU is %trajectories
a sequence of 2D BEV positions for each actor of shape $\mathbb{R}^{N^\mathrm{ref} \times T_f \times 2}$,
and we use finite differences to calculate the heading.
The result is trajectories
$\mathbf{T}^{\mathrm{ref}(i)} \in \mathbb{R}^{N^\mathrm{ref} \times T_f \times 3}$,
describing each objects' BEV pose $\{(x, y, \theta)_{t + s_f}, \dots, (x, y, \theta)_{t + s_f + T_f - 1}\}$.
Note while the GRU sequence length is $T_f + 1$, we only use the last $T_f$ outputs as the future trajectory forecasts.


\subsection{Training}

\paragraph{Augmentations:}

When training \ourmodel{} we adopt standard data augmentations used in prior works~\cite{safdnet,hednet,wang2023dsvt},
including random translations in the x, y, and z directions in the ego coordinate frame with a standard deviation of 0.5m,
random rotation around the z-axis (in the BEV plane) uniformly sampled from $[-\pi/4, \pi/4]$,
random scaling sampled uniformly from $[0.95, 1.05]$,
and flipping with $\{\text{no flip}, \text{x flip}, \text{y flip}, \text{x = y flip}\}$
occurring with equal probability.
We do not use any ground-truth sampling~\cite{yan2018second} because
it is difficult to make consistent and realistic over time and thus incompatible with memory training.


\paragraph{Loss Functions:}

Recall our loss function described in \cref{sec:method-training},
\begin{footnotesize}
\begin{align}
L = L_{\mathrm{rescore}}(\mathbf{C}^\mathrm{merge}) + \sum_{i=1}^I L_{\mathrm{det}}(\mathbf{B}^{\mathrm{ref}(i)}, \mathbf{C}^{\mathrm{ref}(i)}) + L_{\mathrm{for}}(\mathbf{T}^{\mathrm{ref}(i)}),
\end{align}
\end{footnotesize}
which is a combination of
a rescoring loss $L_\mathrm{rescore}$,
a detection refinement loss $L_\mathrm{det}$,
and a forecasting refinement loss $L_\mathrm{for}$, where the detection and forecasting losses are computed at every refinement block.
For brevity, in the above formula we omit the labels in the inputs to the above loss functions.
The labels consist of the bounding boxes for $L_\mathrm{rescore}$ and $L_\mathrm{det}$,
and the ground truth future trajectories for $L_\mathrm{for}$.
% \sergio{I think this last sentence is a bit confusing. Should we say we just omit the labels instead of the inputs?}

Our detection refinement loss $L_\textrm{det}(\mathbf{B}^{\mathrm{ref}(i)}, \mathbf{C}^{\mathrm{ref}(i)})$,
includes a binary focal loss $L_\textrm{det}^\textrm{cls}$ for classification,
and an L1 loss $L_\textrm{det}^\textrm{L1}$ and IoU loss $L_\textrm{det}^\textrm{IoU}$ to regress
the bounding box parameters: $L_\textrm{det} = \lambda_\mathrm{cls} L_\textrm{det}^\textrm{cls} + \lambda_\mathrm{L1} L_\textrm{det}^\textrm{L1} + \lambda_\mathrm{IoU} L_\textrm{det}^\textrm{IoU}$.
For $L_\textrm{det}^\textrm{cls}$ we use focal loss parameters $\alpha = 0.5, \gamma = 2.0$.
We use loss weightings of $\lambda_\mathrm{cls} = 1.0$, $\lambda_\mathrm{L1} = 0.1$, and $\lambda_\mathrm{IoU} = 4.0$.
To calculate the targets for these losses, we first match the detections to the ground truth bounding boxes through bipartite matching as proposed in DETR~\cite{carion2020end}.
We use the same costs for bipartite matching as in DETR except that we omit the L1 component of the box cost, and use 3D IoU instead of 2D IoU.

The rescoring loss $L_\mathrm{rescore}$ is the same as the detection refinement loss described above but without regression $\lambda_{\mathrm{L1}} = \lambda_{\mathrm{IoU}} = 0$.

\subsection{3D Detector Implementation Details}

We use the official implementations for all 3D detectors in our experiments:
single stage CenterPoint~\cite{centerpoint}\footnote{\href{https://github.com/tianweiy/CenterPoint}{https://github.com/tianweiy/CenterPoint}},
HEDNet~\cite{hednet} and SAFDNet~\cite{safdnet}\footnote{\href{https://github.com/zhanggang001/HEDNet}{https://github.com/zhanggang001/HEDNet}},
FCOS3D~\cite{wang2021fcos3d}\footnote{\href{https://github.com/jjw-DL/mmdetection3d\_Noted/tree/master/configs/fcos3d}{https://github.com/jjw-DL/mmdetection3d\_Noted/}},
and BEVMap~\cite{chang2024bevmap}\footnote{\href{https://github.com/mincheoree/BEVMap}{https://github.com/mincheoree/BEVMap}}.

For the LiDAR based models (CenterPoint~\cite{centerpoint}, HEDNet~\cite{hednet}, SAFDNet~\cite{safdnet}),
we train with their official configurations for the WOD.
Neither FOCS3D nor BEVMap provide official configurations for training on AV2;
we train with the ``front ring camera" on AV2 for 3 epochs with a batch size of 16, using the proposed learning rate scheduler.

To obtain $\mathbf{Q}^\mathrm{det}$ from the 3D detectors,
for each object proposal $(x, y, z, l, w, h, \theta)$ in $\mathbf{B} \in \mathbb{R}^{N \times 7}$,
we interpolate the features map directly before the detection header at the projected object centroid $(x, y, z)$.
Concretely, for detectors that produce BEV feature maps before their header
(e.g., CenterPoint~\cite{centerpoint}, HEDNet~\cite{hednet}, SAFDNet~\cite{safdnet}, BEVMap~\cite{chang2024bevmap}),
we perform bi-linear interpolation of the feature maps at the BEV object centroid $(x, y)$.
For detectors that use image-view feature maps, we perform bi-linear
interpolation at the object centroid $(x, y, z)$ projected onto the image plane.
While not experimented with in this work, we could similarly use tri-linear interpolation for detectors that use 3D feature maps.